\documentclass[a4paper,landscape,8pt]{extarticle}
\usepackage[landscape,margin=0.7cm,top=0.5cm,bottom=0.7cm]{geometry}
\usepackage[utf8]{inputenc}
\usepackage[T1]{fontenc}
\usepackage[ngerman]{babel}
\usepackage{helvet}
\renewcommand{\familydefault}{\sfdefault}
\usepackage{xcolor}
\usepackage{tikz}
\usetikzlibrary{positioning, shapes.geometric}
\usepackage{enumitem}
\usepackage{multicol}
\usepackage{array}
\usepackage{fancyhdr}
\usepackage{amssymb}

% Farben
\definecolor{wlanblue}{HTML}{007AFF}
\definecolor{mobilegreen}{HTML}{34C759}
\definecolor{roamingorange}{HTML}{FF9500}
\definecolor{dataviolet}{HTML}{AF52DE}
\definecolor{lightgray}{HTML}{F5F5F5}
\definecolor{bordergray}{HTML}{BBBBBB}
\definecolor{warnred}{HTML}{C62828}

\pagestyle{fancy}
\fancyhf{}
\renewcommand{\headrulewidth}{0pt}
\fancyfoot[C]{\footnotesize Seite \thepage{} von 3}

\setlength{\parindent}{0pt}
\setlength{\parskip}{2pt}

% Kompakte Listen
\setlist{nosep, leftmargin=1em, itemsep=0pt, topsep=1pt}

\begin{document}

% ============================================================================
% SEITE 1: INFORMATIONSBLATT
% ============================================================================

\noindent{\Large\bfseries WLAN \& Mobile Daten: So kommt dein Gerät ins Internet} \hfill
{\footnotesize\textbf{Lernziele:} WLAN vs. Mobilfunk verstehen | Verbinden | Datenverbrauch}

\vspace{1mm}
\hrule
\vspace{1mm}

% Drei-Spalten-Layout
\begin{minipage}[t]{0.32\linewidth}
\begin{center}
\colorbox{wlanblue!10}{%
\begin{minipage}{0.95\linewidth}
\vspace{2mm}
\begin{center}
{\Large\bfseries\textcolor{wlanblue}{WLAN}}\\[1pt]
{\footnotesize\itshape Das Funknetz zu Hause}
\end{center}
\vspace{2mm}
\end{minipage}%
}
\end{center}

\vspace{1mm}

\textbf{Was ist WLAN?} Ein \textbf{kabelloses Netzwerk} -- meist vom Router zu Hause.

\textbf{Analogie:} Wie ein \textbf{unsichtbares Kabel} zwischen deinem Gerät und dem Internet-Anschluss.

\textbf{Vorteile:}
\begin{itemize}
\item \textbf{Kostenlos} (du zahlst nur den Internetanschluss)
\item \textbf{Schnell} (meist schneller als Mobilfunk)
\item \textbf{Unbegrenzt} (kein Datenvolumen)
\end{itemize}

\textbf{Nachteile:}
\begin{itemize}
\item Nur dort, wo ein WLAN ist
\item Braucht \textbf{Passwort} (außer offene Netze)
\end{itemize}

\textbf{WLAN-Symbol:} \raisebox{-0.5ex}{\tikz{\draw[wlanblue, line width=0.5pt] (0,0) arc[start angle=180, end angle=0, radius=0.15] (0,0.1) arc[start angle=180, end angle=0, radius=0.1] (0,0.2) arc[start angle=180, end angle=0, radius=0.05];}} in der Statusleiste

\vspace{2mm}

\textbf{Mit WLAN verbinden:}
\begin{enumerate}
\item Einstellungen → WLAN
\item WLAN \textbf{einschalten} (grüner Schalter)
\item Netzwerk aus Liste wählen
\item \textbf{Passwort} eingeben
\item Fertig -- Gerät merkt sich das Netz!
\end{enumerate}

\end{minipage}%
\hfill%
\begin{minipage}[t]{0.32\linewidth}
\begin{center}
\colorbox{mobilegreen!10}{%
\begin{minipage}{0.95\linewidth}
\vspace{2mm}
\begin{center}
{\Large\bfseries\textcolor{mobilegreen}{Mobile Daten}}\\[1pt]
{\footnotesize\itshape Unterwegs im Internet}
\end{center}
\vspace{2mm}
\end{minipage}%
}
\end{center}

\vspace{1mm}

\textbf{Was sind Mobile Daten?} Internet über das \textbf{Mobilfunknetz} -- wie Telefonieren, aber für Daten.

\textbf{Analogie:} Wie \textbf{Telefonieren} -- überall möglich, wo Empfang ist, aber es kostet.

\textbf{Vorteile:}
\begin{itemize}
\item \textbf{Überall} verfügbar (mit Empfang)
\item Kein Passwort nötig
\item Automatisch verbunden
\end{itemize}

\textbf{Nachteile:}
\begin{itemize}
\item \textbf{Begrenzt} (Datenvolumen im Vertrag)
\item Kann \textbf{teuer} werden (Ausland!)
\item Manchmal langsamer als WLAN
\end{itemize}

\textbf{Symbole in der Statusleiste:}
\begin{itemize}
\item \textbf{5G} = sehr schnell (neueste Technik)
\item \textbf{LTE / 4G} = schnell
\item \textbf{3G} = langsamer
\item \textbf{E} = sehr langsam (Edge)
\end{itemize}

\textbf{Mobile Daten ein/aus:}
\begin{itemize}
\item Einstellungen → Mobilfunk → Mobile Daten
\item Oder: Kontrollzentrum (grüne Antenne)
\end{itemize}

\end{minipage}%
\hfill%
\begin{minipage}[t]{0.32\linewidth}
\begin{center}
\colorbox{lightgray}{%
\begin{minipage}{0.95\linewidth}
\vspace{2mm}
\begin{center}
{\Large\bfseries Der Unterschied}\\[1pt]
{\footnotesize\itshape Wann nutze ich was?}
\end{center}
\vspace{2mm}
\end{minipage}%
}
\end{center}

\vspace{1mm}

\textbf{Automatische Auswahl:}\\
Dein iPhone wählt \textbf{automatisch WLAN}, wenn verfügbar. Nur ohne WLAN werden Mobile Daten genutzt.

\vspace{2mm}

\renewcommand{\arraystretch}{1.4}
\begin{tabular}{@{}p{2cm}p{2.2cm}p{2.2cm}@{}}
& \textbf{\textcolor{wlanblue}{WLAN}} & \textbf{\textcolor{mobilegreen}{Mobil}} \\
\hline
Kosten & Keine & Datenvolumen \\
Reichweite & Nur lokal & Überall \\
Geschw. & Sehr schnell & Schnell \\
Passwort & Ja & Nein \\
\end{tabular}
\renewcommand{\arraystretch}{1}

\vspace{3mm}

\textbf{Faustregel:}
\begin{itemize}
\item \textbf{Zu Hause:} WLAN nutzen
\item \textbf{Unterwegs:} Mobile Daten
\item \textbf{Videos/Updates:} Nur im WLAN!
\end{itemize}

\vspace{2mm}

\begin{center}
\fcolorbox{warnred}{warnred!8}{%
\begin{minipage}{0.92\linewidth}
\vspace{2mm}
\textbf{Achtung:} Große Downloads (Apps, Videos) verbrauchen viel Datenvolumen. Am besten nur im \textbf{WLAN} herunterladen!
\vspace{2mm}
\end{minipage}}
\end{center}

\end{minipage}

\vspace{1mm}

% Zusammenfassung
\begin{center}
\fcolorbox{bordergray}{lightgray}{%
\begin{minipage}{0.98\linewidth}
\centering
\vspace{1.5mm}
\textbf{Merke:}
\textcolor{wlanblue}{\textbf{WLAN}} = zu Hause, kostenlos, unbegrenzt ~$|$~
\textcolor{mobilegreen}{\textbf{Mobile Daten}} = unterwegs, vom Datenvolumen ~$|$~
Dein Gerät bevorzugt automatisch WLAN
\vspace{1.5mm}
\end{minipage}%
}
\end{center}

\newpage

% ============================================================================
% SEITE 2: DATENVOLUMEN & ROAMING
% ============================================================================

\begin{center}
{\LARGE\bfseries Datenvolumen, Roaming \& Einstellungen}\\[3pt]
{\normalsize Was verbraucht Daten? Was kostet im Ausland?}
\end{center}

\vspace{1mm}
\hrule
\vspace{2mm}

\begin{multicols}{2}

\section*{\textcolor{dataviolet}{Datenvolumen verstehen}}

\textbf{Was ist Datenvolumen?} Die \textbf{Menge an Daten}, die du pro Monat über Mobilfunk nutzen kannst.

\textbf{Analogie:} Wie ein \textbf{Wassertank} -- wenn er leer ist, fließt nichts mehr (oder es wird teuer).

\textbf{Typische Verträge:}
\begin{itemize}
\item 1-3 GB = Wenig (nur Nachrichten, E-Mail)
\item 5-10 GB = Mittel (normal surfen)
\item 15+ GB = Viel (auch Videos unterwegs)
\item Unbegrenzt = Flatrate (selten, teuer)
\end{itemize}

\vspace{2mm}

\textbf{Was verbraucht wie viel?}

\renewcommand{\arraystretch}{1.3}
\begin{tabular}{@{}p{4cm}p{3cm}@{}}
WhatsApp-Nachricht & fast nichts \\
E-Mail (ohne Anhang) & ca. 50 KB \\
Foto senden & 1--5 MB \\
1 Stunde Musik streamen & ca. 100 MB \\
1 Stunde Video (Netflix) & ca. 1--3 GB! \\
App-Update & 50--500 MB \\
\end{tabular}
\renewcommand{\arraystretch}{1}

\vspace{2mm}

\textbf{Verbrauch prüfen:}\\
Einstellungen → Mobilfunk → \glqq Aktueller Zeitraum\grqq{}

Dort siehst du, \textbf{welche App} wie viel verbraucht hat.

\vspace{2mm}

\textbf{Daten sparen:}
\begin{itemize}
\item Videos nur im WLAN schauen
\item App-Updates nur im WLAN (Einstellungen → App Store)
\item \glqq Datensparmodus\grqq{} aktivieren
\end{itemize}

\columnbreak

\section*{\textcolor{roamingorange}{Roaming -- Im Ausland}}

\textbf{Was ist Roaming?} Dein Handy nutzt ein \textbf{fremdes Netz} im Ausland.

\textbf{Analogie:} Wie \textbf{Telefonieren mit Hoteltelefon} -- funktioniert, aber das Hotel berechnet extra.

\vspace{2mm}

\textbf{In der EU:}
\begin{itemize}
\item \textbf{Keine Extrakosten!} (seit 2017)
\item Du nutzt dein deutsches Datenvolumen
\item Telefonieren wie zu Hause
\end{itemize}

\textbf{Außerhalb der EU (z.B. Schweiz, USA, Türkei):}
\begin{itemize}
\item \textbf{Sehr teuer!} (mehrere Euro pro MB)
\item Mobile Daten \textbf{unbedingt ausschalten}
\item Nur WLAN nutzen
\end{itemize}

\vspace{2mm}

\begin{center}
\fcolorbox{warnred}{warnred!8}{%
\begin{minipage}{0.9\linewidth}
\vspace{2mm}
\textbf{Vor der Reise (außerhalb EU):}\\[2pt]
Einstellungen → Mobilfunk → \textbf{Datenroaming AUS}\\[2pt]
Sonst drohen Rechnungen von \textbf{hunderten Euro}!
\vspace{2mm}
\end{minipage}}
\end{center}

\vspace{3mm}

\section*{Wichtige Einstellungen}

\textbf{Einstellungen → WLAN:}
\begin{itemize}
\item WLAN ein/aus
\item Bekannte Netzwerke verwalten
\item \glqq Automatisch verbinden\grqq{}
\end{itemize}

\textbf{Einstellungen → Mobilfunk:}
\begin{itemize}
\item Mobile Daten ein/aus
\item Datenroaming ein/aus
\item Pro App: Daten erlauben ja/nein
\end{itemize}

\vspace{2mm}

\begin{center}
\fcolorbox{wlanblue}{wlanblue!8}{%
\begin{minipage}{0.9\linewidth}
\vspace{2mm}
\textbf{Tipp: WLAN-Unterstützung}\\[2pt]
Einstellungen → Mobilfunk → ganz unten: \glqq WLAN-Unterstützung\grqq{} \textbf{ausschalten}. Sonst springt das iPhone bei schwachem WLAN auf Mobile Daten um.
\vspace{2mm}
\end{minipage}}
\end{center}

\end{multicols}

\newpage

% ============================================================================
% SEITE 3: ÜBUNGEN & CHECKLISTE
% ============================================================================

\begin{center}
{\LARGE\bfseries Übungen \& Checkliste}\\[3pt]
{\normalsize Teste dein Wissen über WLAN und Mobile Daten}
\end{center}

\vspace{1mm}
\hrule
\vspace{2mm}

\begin{multicols}{2}

\textbf{\large Teil A: Richtig oder Falsch?}

\vspace{2mm}

\renewcommand{\arraystretch}{1.4}
\begin{tabular}{@{}p{8cm}cc@{}}
& R & F \\
1. WLAN ist immer kostenlos. & $\square$ & $\square$ \\
2. Mobile Daten funktionieren nur zu Hause. & $\square$ & $\square$ \\
3. Videos verbrauchen viel Datenvolumen. & $\square$ & $\square$ \\
4. In der EU kostet Roaming extra. & $\square$ & $\square$ \\
5. Das iPhone bevorzugt automatisch WLAN. & $\square$ & $\square$ \\
6. WhatsApp-Nachrichten verbrauchen viel Daten. & $\square$ & $\square$ \\
\end{tabular}
\renewcommand{\arraystretch}{1}

\textit{\footnotesize Lösungen: R, F, R, F, R, F}

\vspace{4mm}

\textbf{\large Teil B: Was verbraucht mehr?}

Kreuze an, was \textbf{mehr Datenvolumen} verbraucht:

\vspace{2mm}

\begin{tabular}{@{}p{5cm}p{5cm}@{}}
$\square$ E-Mail lesen & $\square$ Video schauen \\[3mm]
$\square$ WhatsApp-Text & $\square$ WhatsApp-Foto \\[3mm]
$\square$ App herunterladen & $\square$ App öffnen \\[3mm]
$\square$ Musik streamen & $\square$ Wetter-App prüfen \\
\end{tabular}

\vspace{4mm}

\textbf{\large Teil C: Praktische Übungen}

Führe diese Aktionen auf deinem iPhone aus:

\vspace{2mm}

\begin{itemize}[label=$\square$, itemsep=3pt]
\item Öffne \textbf{Einstellungen → WLAN}
\item Prüfe, mit welchem Netzwerk du verbunden bist
\item Öffne \textbf{Einstellungen → Mobilfunk}
\item Finde heraus, wie viel Daten du verbraucht hast
\item Prüfe, ob \textbf{Datenroaming} ein oder aus ist
\item Öffne das \textbf{Kontrollzentrum} und finde WLAN/Mobile Daten
\end{itemize}

\columnbreak

\textbf{\large Teil D: Situationen}

Was machst du?

\vspace{2mm}

\textbf{Situation 1:} Du bist im Café und willst ein Video schauen.

\rule{7cm}{0.4pt}

\vspace{3mm}

\textbf{Situation 2:} Du fliegst morgen in die Türkei (kein EU-Land).

\rule{7cm}{0.4pt}

\vspace{4mm}

\textbf{\large Checkliste: Das weiß ich jetzt}

\vspace{2mm}

\begin{center}
\fcolorbox{wlanblue}{wlanblue!10}{%
\begin{minipage}{0.95\linewidth}
\vspace{2mm}
\textbf{\textcolor{wlanblue}{WLAN}}
\begin{itemize}[label=$\square$, itemsep=2pt]
\item Ich weiß, was WLAN ist und wie ich mich verbinde
\item Ich weiß, dass WLAN \textbf{unbegrenzt} und meist kostenlos ist
\item Ich kann das WLAN-Symbol erkennen
\end{itemize}
\vspace{2mm}
\end{minipage}}
\end{center}

\vspace{2mm}

\begin{center}
\fcolorbox{mobilegreen}{mobilegreen!10}{%
\begin{minipage}{0.95\linewidth}
\vspace{2mm}
\textbf{\textcolor{mobilegreen}{Mobile Daten}}
\begin{itemize}[label=$\square$, itemsep=2pt]
\item Ich weiß, dass Mobile Daten vom \textbf{Datenvolumen} abgehen
\item Ich weiß, was viel Daten verbraucht (Videos!)
\item Ich kann meinen Verbrauch prüfen
\end{itemize}
\vspace{2mm}
\end{minipage}}
\end{center}

\vspace{2mm}

\begin{center}
\fcolorbox{roamingorange}{roamingorange!10}{%
\begin{minipage}{0.95\linewidth}
\vspace{2mm}
\textbf{\textcolor{roamingorange}{Roaming}}
\begin{itemize}[label=$\square$, itemsep=2pt]
\item Ich weiß, dass in der \textbf{EU keine Extrakosten} anfallen
\item Ich weiß, dass \textbf{außerhalb der EU} Roaming sehr teuer ist
\item Ich kann Datenroaming \textbf{ausschalten}
\end{itemize}
\vspace{2mm}
\end{minipage}}
\end{center}

\vspace{2mm}

\begin{center}
\fcolorbox{bordergray}{white}{%
\begin{minipage}{0.95\linewidth}
\vspace{2mm}
\textbf{Meine Notizen}\\[2pt]
Mein WLAN-Name zu Hause: \dotfill\\[2mm]
Mein Datenvolumen pro Monat: \dotfill\\[2mm]
Mein Mobilfunkanbieter: \dotfill
\vspace{2mm}
\end{minipage}}
\end{center}

\end{multicols}

\vfill

\begin{center}
\footnotesize\textit{Stand: \today{} -- Zu Hause WLAN, unterwegs Mobile Daten. Im Nicht-EU-Ausland: Roaming ausschalten!}
\end{center}

\end{document}
